% =====================================================================
% ADDITION FOR 06-hardware.tex  (prepared for Rev J)
% Insert in \section{Servers, Computers, and Operator Stations},
% directly AFTER the "Industrial Keyboards" subsubsection (and its
% italic notes paragraph) and BEFORE
% \subsection{Class 1 -- Div 2 Zone Hardware}\label{sec:c1d2}
% =====================================================================

\subsection{Operator Station Mounting Height}\label{sec:hmi-height}
These rules apply to every operator display: panel-door HMIs and touch monitors, swing-arm monitors, wall-mounted plant-floor stations, and control-room desks. Measure every height from the surface the operator actually stands or sits on (finished floor, platform, or grating), not from the floor at the base of the panel.

\begin{itemize}[itemsep=2pt]
  \item \textbf{Standing operator:} the screen centerline must be \textbf{1500~mm (59~in) $\pm$100~mm} above the standing surface.
  \item \textbf{Whole active area:} every touch target and on-screen control must sit between \textbf{600~mm (24~in)} and \textbf{1700~mm (67~in)} above the standing surface.
  \item \textbf{Seated operator} (control-room desk or console): the screen centerline should be about 1050--1200~mm (41--47~in) above the floor, with the top of the screen at or below seated eye level.
\end{itemize}

\textbf{Why 1500~mm.} A screen is easiest to read when it sits at or slightly below eye level, so the operator looks about 10--20\textdegree\ downward.\footnote{ISO 9241-5, \emph{Ergonomic requirements for office work with visual display terminals --- Part 5: Workstation layout and postural requirements}; ANSI/HFES 100-2007, \emph{Human Factors Engineering of Computer Workstations}.} Standing eye height ranges from roughly 1450~mm (a small woman) to roughly 1750~mm (a tall man). A 1500~mm centerline keeps the screen readable and the touch targets within easy reach across that range, without anyone stretching up or stooping down.

\textbf{Why 600~mm and 1700~mm.} A touch HMI is an operator control device. NFPA 79 and IEC 60204-1 require hand-operated control devices to be at least 0.6~m above the servicing level and readily within reach of the normal working position.\footnote{NFPA 79-2024, Chapter 10 (operator interface and control devices); IEC 60204-1:2016+A1:2021, 10.1.2.} The 1700~mm upper limit keeps the top of the screen within comfortable reach for a short operator.\footnote{ISO 14738:2002, \emph{Safety of machinery --- Anthropometric requirements for the design of workstations at machinery}.}

\textbf{Accessibility.} If the project or site requires ADA accessibility at an operator station, the touch area must instead fall within the 380--1220~mm (15--48~in) reach range, which puts the centerline at about 1100--1200~mm. The project specification must say whether this applies.\footnote{2010 ADA Standards for Accessible Design, 308 (reach ranges).}

\textbf{Practical points:}
\begin{itemize}[itemsep=2pt]
  \item Mount the HMI on the enclosure door or a swing arm, not deep inside the panel, so the operator can reach it squarely.
  \item Avoid placing the screen directly under overhead lights or facing a window, where glare washes it out.
  \item If the screen has to sit low, tilt it slightly upward toward the operator.
  \item The panel layout drawing (Section~\ref{sec:drawing-package}) must dimension the HMI centerline from the operator's standing surface.
\end{itemize}

\note{The 1700~mm upper limit is an EFC ergonomic choice based on ISO 14738 rather than a hard code limit. Confirm it during review, along with whether any sites require ADA-compliant operator stations.}
